\documentclass[11pt, a4paper]{article}

% Gemeinsame Style-Definitionen (TEXINPUTS muss templates/ enthalten — siehe scripts/build_tex.py)
% bfw-style.tex — gemeinsame Style-Definitionen für alle Handouts/Aufgabenhefte
% Voraussetzung: Kompilation mit xelatex (wegen fontspec)

\usepackage[ngerman]{babel}
\usepackage{geometry}
\geometry{a4paper, top=2.2cm, bottom=2.2cm, left=2.3cm, right=2.3cm}

\usepackage{fontspec}
% Fallback-Kette: Source Sans Pro -> Calibri -> Segoe UI -> Latin Modern Sans
% (Latin Modern ist immer mit MiKTeX vorhanden)
\IfFontExistsTF{Source Sans Pro}{%
  \setmainfont{Source Sans Pro}[Ligatures=TeX]%
}{\IfFontExistsTF{Calibri}{%
  \setmainfont{Calibri}[Ligatures=TeX]%
}{\IfFontExistsTF{Segoe UI}{%
  \setmainfont{Segoe UI}[Ligatures=TeX]%
}{%
  \setmainfont{Latin Modern Sans}[Ligatures=TeX]%
}}}
\IfFontExistsTF{Source Code Pro}{%
  \setmonofont{Source Code Pro}[Scale=0.92]%
}{\IfFontExistsTF{Consolas}{%
  \setmonofont{Consolas}[Scale=0.92]%
}{%
  \setmonofont{Latin Modern Mono}[Scale=0.92]%
}}

\usepackage{xcolor}
% Farbpalette: hell, freundlich, modern
\definecolor{bfwPrimary}{HTML}{0EA5A4}    % Petrol/Türkis - Primär
\definecolor{bfwPrimaryDark}{HTML}{0F766E} % dunkler Petrol für Headings
\definecolor{bfwAccent}{HTML}{F59E0B}     % Amber - Akzent
\definecolor{bfwSuccess}{HTML}{10B981}    % Grün - Erfolg / Lösung
\definecolor{bfwWarn}{HTML}{F97316}       % Orange - Achtung
\definecolor{bfwInk}{HTML}{0F172A}        % fast Schwarz
\definecolor{bfwInkSoft}{HTML}{475569}    % gedämpftes Grau
\definecolor{bfwBgSoft}{HTML}{F8FAFC}     % off-white
\definecolor{bfwBgTeal}{HTML}{ECFEFF}     % helles Türkis
\definecolor{bfwBgAmber}{HTML}{FEF3C7}    % helles Gelb
\definecolor{bfwBgRose}{HTML}{FFE4E6}     % helles Rosé für Eselsbrücken

\usepackage[colorlinks=true, linkcolor=bfwPrimaryDark, urlcolor=bfwPrimaryDark]{hyperref}
\usepackage{enumitem}
\setlist[itemize]{leftmargin=*, itemsep=2pt, topsep=2pt}
\setlist[enumerate]{leftmargin=*, itemsep=2pt, topsep=2pt}

\usepackage[most]{tcolorbox}
\usepackage{booktabs}
\usepackage{tikz}
\usetikzlibrary{decorations.pathmorphing, positioning, shapes.geometric, arrows.meta}

% Merkkästchen — wichtige Definition / Kernpunkt
\newtcolorbox{merksatz}[1][]{%
  enhanced, breakable,
  colback=bfwBgTeal,
  colframe=bfwPrimary,
  boxrule=0.6pt,
  arc=4pt,
  left=10pt, right=10pt, top=8pt, bottom=8pt,
  fonttitle=\bfseries\color{bfwPrimaryDark},
  title={\faLightbulb\ Merksatz},
  #1
}

% Eselsbrücke — Mnemoniks
\newtcolorbox{eselsbruecke}[1][]{%
  enhanced, breakable,
  colback=bfwBgRose,
  colframe=bfwAccent,
  boxrule=0.6pt,
  arc=4pt,
  left=10pt, right=10pt, top=8pt, bottom=8pt,
  fonttitle=\bfseries\color{bfwWarn},
  title={Eselsbrücke},
  #1
}

% Beispiel-Box
\newtcolorbox{beispiel}[1][]{%
  enhanced, breakable,
  colback=bfwBgSoft,
  colframe=bfwInkSoft,
  boxrule=0.4pt,
  arc=3pt,
  left=10pt, right=10pt, top=8pt, bottom=8pt,
  fonttitle=\bfseries\color{bfwInk},
  title={Beispiel},
  #1
}

% Lösungs-Box (für Aufgabenhefte)
\newtcolorbox{loesung}[1][]{%
  enhanced, breakable,
  colback=bfwBgAmber!40,
  colframe=bfwSuccess,
  boxrule=0.6pt,
  arc=3pt,
  left=10pt, right=10pt, top=8pt, bottom=8pt,
  fonttitle=\bfseries\color{bfwSuccess!70!black},
  title={Lösung},
  #1
}

% Glossar-Eintrag
\newcommand{\glossareintrag}[2]{%
  \par\noindent\textbf{\color{bfwPrimaryDark}#1} \enspace #2\par\smallskip
}

% Section-Stil — etwas farbiger als Standard
\usepackage{titlesec}
\titleformat{\section}
  {\Large\bfseries\color{bfwPrimaryDark}}
  {\thesection}{0.6em}{}
\titleformat{\subsection}
  {\large\bfseries\color{bfwPrimary}}
  {\thesubsection}{0.5em}{}
\titleformat{\subsubsection}
  {\normalsize\bfseries\color{bfwInk}}
  {\thesubsubsection}{0.4em}{}

% TikZ-Stil "skizzenhaft" — etwas weicher als Rough.js, aber stabil
\tikzset{
  roughbox/.style = {
    draw=bfwPrimaryDark, thick, fill=bfwBgTeal,
    rounded corners=4pt, inner sep=6pt
  },
  roughaccent/.style = {
    draw=bfwAccent, thick, fill=bfwBgAmber,
    rounded corners=4pt, inner sep=6pt
  },
  roughline/.style = {
    draw=bfwInkSoft, thick, -{Stealth[length=2.5mm]}
  }
}

% Bullet-Symbole (bewusst kein FontAwesome — vermeidet Font-Installations-Probleme)
\newcommand{\faLightbulb}{$\bullet$}
\newcommand{\faBookmark}{$\bullet$}

% Header/Footer
\usepackage{fancyhdr}
\pagestyle{fancy}
\fancyhf{}
\renewcommand{\headrulewidth}{0pt}
\fancyfoot[L]{\small\color{bfwInkSoft}\thepage}
\fancyfoot[R]{\small\color{bfwInkSoft} BFW M-064 Beschaffung im Büromanagement}


\title{\vspace*{-1.5cm}\Huge\bfseries\color{bfwPrimaryDark}Session~3: Das Unternehmen --- Firma, Rechtsformen \& Organisation\\[0.4em]
  \large\color{bfwInkSoft}\textnormal{Wer ist Kaufmann, wie heißt die Firma und wie ist der Betrieb aufgebaut?}}
\author{\textsf{Lernfeld 1 --- Die eigene Rolle im Betrieb mitgestalten \textperiodcentered{} \small\copyright\ Can Siebert\,\textperiodcentered\,2026}}
\date{}

\begin{document}
\maketitle
\thispagestyle{empty}

\vspace{1em}
\begin{merksatz}[title={\bfseries Worum geht es in dieser Session?}]
Jedes Unternehmen hat einen rechtlichen Rahmen: Wer als \emph{Kaufmann} gilt, unter
welchem Namen (\emph{Firma}) es auftritt, in welcher \emph{Rechtsform} es geführt wird
und wie es intern \emph{organisiert} ist. Diese Fragen entscheiden über Haftung, Kapital,
Buchführungspflichten und die Verteilung von Verantwortung. In diesem Handout lernen Sie
den Kaufmannsbegriff, die Firma und ihre Grundsätze, das Handelsregister, die wichtigsten
Rechtsformen, die Vollmachten (Prokura und Handlungsvollmacht) sowie die Aufbau- und
Ablauforganisation kennen --- ergänzt um Leitbild, Unternehmensziele und Nachhaltigkeit.
Das ist Prüfungsstoff für die IHK-Abschlussprüfung.
\end{merksatz}

\tableofcontents
\vspace{1em}
\hrule
\vspace{1em}

\section{Der Kaufmannsbegriff nach HGB}

Das Handelsgesetzbuch (HGB) gilt nicht für jeden, der etwas verkauft, sondern für den
\emph{Kaufmann}. Ob jemand Kaufmann ist, entscheidet über Buchführungspflichten, die
Firmenführung und besondere Rechte und Pflichten. Man unterscheidet drei Grundtypen.

\begin{itemize}
  \item \textbf{Istkaufmann} (§ 1 HGB) --- Wer ein Handelsgewerbe betreibt, das nach Art
    und Umfang einen in kaufmännischer Weise eingerichteten Geschäftsbetrieb erfordert,
    ist \emph{kraft Gesetzes} Kaufmann. Er ist zur Eintragung ins Handelsregister
    verpflichtet; diese Eintragung ist aber nur \emph{deklaratorisch} (rechts\-bekundend)
    --- er ist bereits vorher Kaufmann.
  \item \textbf{Kannkaufmann} (§§ 2, 3 HGB) --- Kleingewerbetreibende sowie land- und
    forstwirtschaftliche Unternehmen sind zunächst keine Kaufleute. Sie \emph{können}
    aber auf eigenen Wunsch Kaufmann werden, indem sie sich freiwillig ins Handelsregister
    eintragen lassen.
  \item \textbf{Formkaufmann} (z.\,B.\ § 6 HGB) --- Kapitalgesellschaften wie AG, GmbH,
    KGaA und die eingetragene Genossenschaft (eG) sind allein wegen ihrer \emph{Rechtsform}
    Kaufmann. Sie werden erst mit der Eintragung ins Handelsregister Kaufmann; diese
    Eintragung wirkt \emph{konstitutiv} (rechts\-begründend).
\end{itemize}

\begin{merksatz}
Der Kleingewerbetreibende sollte sich die freiwillige Eintragung als Kannkaufmann gut
überlegen: Mit ihr treffen ihn die \emph{Buchführungspflichten} nach HGB sowie besondere
\emph{Untersuchungs- und Rügepflichten} im Kaufvertrag.
\end{merksatz}

\begin{eselsbruecke}
\textbf{„Ist von selbst, Kann freiwillig, Form durch Rechtsform.”} \emph{Ist}kaufmann ist
Kaufmann von selbst, \emph{Kann}kaufmann wird es freiwillig durch Eintragung, \emph{Form}\-kaufmann
ist es kraft seiner Rechtsform.
\end{eselsbruecke}

\section{Die Firma --- der Name des Kaufmanns}

Die \textbf{Firma} ist nach § 17 HGB der Name, unter dem ein Kaufmann seine Geschäfte
betreibt und die Unterschrift abgibt. Achtung: Umgangssprachlich meint „Firma" das
Unternehmen --- rechtlich ist sie nur der \emph{Handelsname}. Die Firma besteht aus einem
\emph{Firmenkern} (Name, Gegenstand oder Fantasiebezeichnung) und einem \emph{Firmenzusatz}
(vor allem der Rechtsformzusatz, z.\,B.\ e.K., OHG, KG, GmbH, AG).

\subsection{Firmenarten}

\begin{center}
\begin{tabular}{p{3.1cm} p{7.2cm} p{2.6cm}}
\toprule
\textbf{Firmenart} & \textbf{Firmenkern} & \textbf{Beispiel}\\
\midrule
Personenfirma & Namen der beteiligten Personen & Gilles e.K.\\
Sachfirma & Gegenstand des Unternehmens & Balthasar Transporte GmbH\\
Mischfirma & Name \emph{und} Gegenstand & Duisdorfer BüroKonzept KG\\
Fantasiefirma & frei erfundener Fantasiename & Kalyptus OHG\\
\bottomrule
\end{tabular}
\end{center}

\subsection{Die Firmengrundsätze}

Bei der Wahl der Firma müssen fünf \emph{Firmengrundsätze} beachtet werden.

\begin{enumerate}
  \item \textbf{Firmenwahrheit} (§ 18 HGB) --- Die Firma muss zur Kennzeichnung geeignet
    sein und Unterscheidungskraft besitzen; sie darf nicht über geschäftliche Verhältnisse
    irreführen.
  \item \textbf{Firmenklarheit} (§ 18 Abs.\ 2 HGB) --- Die Firma darf keine Angaben
    enthalten, die zur Irreführung geeignet sind; sie muss klar und verständlich sein.
  \item \textbf{Firmenausschließlichkeit} (§ 30 HGB) --- Jede neue Firma muss sich von
    bereits am selben Ort eingetragenen Firmen deutlich unterscheiden; bei Namens\-gleichheit
    ist ein unterscheidender Zusatz nötig.
  \item \textbf{Firmenbeständigkeit} (§§ 22, 23 HGB) --- Wer ein Handelsgeschäft erwirbt,
    darf die bisherige Firma fortführen. Die Firma kann aber nicht \emph{ohne} das
    Handelsgeschäft veräußert werden.
  \item \textbf{Firmenöffentlichkeit} (§ 29 HGB) --- Jeder Kaufmann ist verpflichtet,
    seine Firma, den Ort und die inländische Geschäftsanschrift zur Eintragung ins
    Handelsregister anzumelden.
\end{enumerate}

\begin{beispiel}
Die „Duisdorfer BüroKonzept KG" (stellt Büroausstattung her) dürfte sich nicht
„Duisdorfer Immobilien KG" nennen --- das verstößt gegen die \emph{Firmenwahrheit}, weil
„Immobilien" nicht dem tatsächlichen Gegenstand entspricht und in die Irre führen würde.
\end{beispiel}

\begin{eselsbruecke}
\textbf{„Wahr, Klar, Ausschließlich, Beständig, Öffentlich”} --- die fünf Firmen\-grundsätze
in einem Satz. Merken Sie sich vor allem: Die Firma muss stimmen (\emph{wahr}), verständlich
sein (\emph{klar}) und einmalig sein (\emph{ausschließlich}).
\end{eselsbruecke}

\section{Das Handelsregister}

Das \textbf{Handelsregister} ist ein elektronisch geführtes amtliches Verzeichnis, das
vom Amtsgericht des jeweiligen Bezirks geführt wird (§ 8 HGB). Jeder darf zu
Informationszwecken Einsicht nehmen (§ 9 HGB) --- so können auch Auszubildende
Informationen über ihren Ausbildungsbetrieb erhalten.

\begin{center}
\begin{tabular}{p{3.2cm} p{9.4cm}}
\toprule
\textbf{Abteilung} & \textbf{Eingetragene Unternehmen}\\
\midrule
Abteilung A & Einzelkaufleute (e.K.) und Personengesellschaften (OHG, KG)\\
Abteilung B & Kapitalgesellschaften (GmbH, AG)\\
\bottomrule
\end{tabular}
\end{center}

\medskip
Genossenschaften (eG) werden im gesonderten \emph{Genossenschaftsregister}, Vereine
(e.\,V.) im \emph{Vereinsregister} eingetragen.

\subsection{Konstitutive und deklaratorische Wirkung}

Nicht jede Eintragung hat dieselbe Wirkung:

\begin{itemize}
  \item \textbf{Deklaratorisch (rechtsbekundend)} --- Die Eintragung bestätigt nur eine
    bereits bestehende Rechtslage, z.\,B.\ die Eintragung des Istkaufmanns oder die
    Eintragung/Löschung einer Prokura.
  \item \textbf{Konstitutiv (rechtsbegründend)} --- Die Rechtslage entsteht erst mit der
    Eintragung, z.\,B.\ die Entstehung einer GmbH oder AG.
\end{itemize}

\subsection{Publizität des Handelsregisters (§ 15 HGB)}

Das Handelsregister genießt \emph{öffentlichen Glauben}. Nach § 15 HGB gilt:

\begin{merksatz}
Solange eine eintragungspflichtige Tatsache \emph{nicht} eingetragen und bekanntgemacht
ist, kann sie einem gutgläubigen Dritten \emph{nicht} entgegengehalten werden. Ist sie
dagegen eingetragen und bekanntgemacht, muss ein Dritter sie \emph{gegen sich gelten
lassen}. Jeder darf also auf den Registerinhalt vertrauen.
\end{merksatz}

\begin{beispiel}
Die Prokura von Frau Sommer wurde widerrufen; der Widerruf ist ordnungsgemäß eingetragen
und bekanntgemacht. Legt sie dennoch eine alte Vollmachtsurkunde vor und bestellt Waren,
muss der Geschäftspartner den eingetragenen Widerruf gegen sich gelten lassen --- das
Unternehmen ist nicht gebunden. Der Partner muss sich direkt an Frau Sommer halten.
\end{beispiel}

\section{Rechtsformen der Unternehmen}

Die Rechtsform bestimmt vor allem \emph{Haftung, Kapital, Geschäftsführung, Vertretung}
und die \emph{Organe}. Man unterscheidet Einzelunternehmung, Personengesellschaften und
Kapitalgesellschaften.

\subsection{Einzelunternehmung und Personengesellschaften}

Die \textbf{Einzelunternehmung} (e.K., e.\,Kfm., e.\,Kffr.) wird von einer Person
gegründet, die das Kapital aus ihrem Privatvermögen aufbringt, allein die Geschäfte führt
und das Unternehmen nach außen vertritt. Der Gewinn steht ihr allein zu --- sie haftet
aber auch \emph{unbeschränkt mit ihrem gesamten Privatvermögen}.

Die \textbf{Gesellschaft bürgerlichen Rechts (GbR)} ist die einfachste Personen\-gesellschaft
für nicht-kaufmännische Zwecke (z.\,B.\ Freiberufler, Gelegenheitsgesellschaften). Auch
hier haften die Gesellschafter persönlich und unbeschränkt. Betreibt eine GbR ein
Handelsgewerbe, wird sie zur OHG.

Die \textbf{offene Handelsgesellschaft (OHG)} wird von zwei oder mehr Personen gegründet.
Jeder Gesellschafter ist allein zur Geschäftsführung und zur Vertretung berechtigt. Die
Gesellschafter haften \emph{unbeschränkt, unmittelbar und solidarisch} (gesamt\-schuld\-nerisch):
Gläubiger können sich an einen einzelnen zahlungsfähigen Gesellschafter wenden und die
gesamte Forderung eintreiben.

Die \textbf{Kommanditgesellschaft (KG)} hat mindestens einen \emph{Komplementär}
(Vollhafter) und mindestens einen \emph{Kommanditisten} (Teilhafter). Nur die
Komplementäre führen die Geschäfte und vertreten die KG; sie haften unbeschränkt.
Kommanditisten haften nur \emph{bis zur Höhe ihrer Kapitaleinlage}.

\begin{center}
\begin{tabular}{p{2.6cm} p{3.3cm} p{3.3cm} p{3.3cm}}
\toprule
\textbf{Merkmal} & \textbf{Einzelunt.\ (e.K.)} & \textbf{OHG} & \textbf{KG}\\
\midrule
Gründung & 1 Person & $\geq$ 2 Personen & $\geq$ 1 Komplementär + 1 Kommanditist\\
Kapital & Privatvermögen & Privatvermögen der Gesellschafter & Privatvermögen\\
Geschäftsführung & Einzelunternehmer & jeder Gesellschafter allein & nur Komplementäre\\
Vertretung & Einzelunt., Prokurist & jeder Gesellschafter, Prokurist & nur Komplementäre, Prokurist\\
Haftung & unbeschränkt, privat & unbeschränkt, unmittelbar, solidarisch & Komplementär unbeschränkt; Kommanditist bis Einlage\\
\bottomrule
\end{tabular}
\end{center}

\subsection{Kapitalgesellschaften und Genossenschaft}

Kapitalgesellschaften sind \emph{juristische Personen} mit eigener Rechtspersönlichkeit;
die Haftung ist grundsätzlich auf das Gesellschaftsvermögen beschränkt.

\begin{itemize}
  \item \textbf{GmbH} --- Mindeststammkapital \textbf{25.000,00 €}. Organe:
    \emph{Gesellschafterversammlung} (bestellt die Geschäftsführer) und ein oder mehrere
    \emph{Geschäftsführer}. Vertretung durch Geschäftsführer oder Prokuristen. Haftung
    beschränkt auf das GmbH-Vermögen; solange Einlagen nicht voll eingezahlt sind, besteht
    eine \emph{Nachschusspflicht} bis zur Höhe des Geschäftsanteils.
  \item \textbf{UG (haftungsbeschränkt)} --- die „Mini-GmbH", gründbar ab
    \textbf{1,00 €} Stammkapital. Gleiche beschränkte Haftung wie die GmbH, aber Gewinne
    dürfen nicht voll ausgeschüttet werden, bis 25.000 € Stammkapital angespart sind.
  \item \textbf{AG} --- Grundkapital mindestens \textbf{50.000,00 €}, zerlegt in Aktien.
    Organe: \emph{Hauptversammlung} (Aktionäre), \emph{Vorstand} (Geschäftsführung) und
    \emph{Aufsichtsrat} (Kontrolle). Vertretung durch Vorstand oder Prokuristen;
    Gewinn\-ausschüttung an Aktionäre heißt \emph{Dividende}.
  \item \textbf{eG (eingetragene Genossenschaft)} --- juristische Person, der \emph{Förderung
    ihrer Mitglieder} verpflichtet. Kapital aus Genossenschaftsanteilen der Mitglieder.
    Organe: \emph{General-/Vertreterversammlung}, gewählter \emph{Vorstand} und
    \emph{Aufsichtsrat}. Haftung mit dem Gesellschaftsvermögen.
\end{itemize}

\begin{center}
\begin{tabular}{p{2.4cm} p{3.4cm} p{3.4cm} p{3.4cm}}
\toprule
\textbf{Merkmal} & \textbf{GmbH} & \textbf{AG} & \textbf{eG}\\
\midrule
Kapital & mind.\ 25.000 € Stammkapital & mind.\ 50.000 € Grundkapital & Genossenschafts\-anteile\\
Eigentümer & Gesellschafter & Aktionäre & Mitglieder\\
Geschäftsführung & bestellter Geschäftsführer & bestellter Vorstand & gewählter Vorstand\\
Vertretung & Geschäftsführer, Prokurist & Vorstand, Prokurist & Vorstand, Prokurist\\
Haftung & Gesellschafts\-vermögen & Gesellschafts\-vermögen & Gesellschafts\-vermögen\\
\bottomrule
\end{tabular}
\end{center}

\begin{eselsbruecke}
\textbf{„Persönlich haftet man privat, Kapitalgesellschaft nur mit Kapital.”} Bei
Einzelunternehmen, GbR und OHG haftet man mit dem \emph{Privatvermögen}; bei GmbH, UG, AG
und eG ist die Haftung auf das \emph{Gesellschaftsvermögen} beschränkt.
\end{eselsbruecke}

\section{Vollmachten: Prokura und Handlungsvollmacht}

Kaufleute können sich im Außenverhältnis vertreten lassen. Die beiden wichtigsten
kaufmännischen Vollmachten sind die Prokura und die Handlungsvollmacht.

\subsection{Prokura (§§ 48 ff.\ HGB)}

Die \textbf{Prokura} wird nur vom Inhaber (oder seinem gesetzlichen Vertreter) durch
ausdrückliche Erklärung erteilt und muss ins Handelsregister eingetragen werden. Sie
ermächtigt zu \emph{allen} gerichtlichen und außergerichtlichen Geschäften, die der
Betrieb eines Handelsgewerbes mit sich bringt. \textbf{Ausnahme:} Die Veräußerung und
Belastung von Grundstücken ist nur bei besonders erteilter Befugnis erlaubt. Ein
Prokurist unterschreibt mit dem Zusatz \emph{„ppa."} (per procura).

\subsection{Handlungsvollmacht (§ 54 HGB)}

Die \textbf{Handlungsvollmacht} wird formlos erteilt und \emph{nicht} ins Handelsregister
eingetragen. Ihr Umfang ist enger und richtet sich nach dem gewöhnlichen Geschäftsbetrieb.
Man unterscheidet:

\begin{itemize}
  \item \textbf{Generalvollmacht} --- alle gewöhnlichen Geschäfte des Betriebs; Unterschrift
    mit \emph{„i.\,V."} (in Vollmacht).
  \item \textbf{Art-/Gattungsvollmacht} --- eine bestimmte Art von Geschäften (z.\,B.\
    regelmäßiges Kassieren); Unterschrift mit \emph{„i.\,A."} (im Auftrag).
  \item \textbf{Einzel-/Spezialvollmacht} --- ein einzelnes Geschäft (z.\,B.\ einmaliger
    Einkauf); Unterschrift mit \emph{„i.\,A."}.
\end{itemize}

Grundstücksgeschäfte, Wechselverbindlichkeiten, Darlehensaufnahme und Prozessführung
darf ein Handlungsbevollmächtigter nur mit besonderer Befugnis vornehmen.

\begin{beispiel}
Auszubildende der Duisdorfer BüroKonzept KG besitzen häufig eine \emph{Artvollmacht}
(z.\,B.\ regelmäßiges Kassieren) und/oder eine \emph{Einzelvollmacht} (z.\,B.\ einmaliger
Kauf von Briefmarken). Mit fortschreitender Ausbildung wächst der Handlungsrahmen.
\end{beispiel}

\section{Aufbau- und Ablauforganisation}

Die \textbf{Aufbauorganisation} legt das „Gerüst" des Betriebs fest: welche Stellen es
gibt und wie sie zueinander stehen. Grundbegriffe:

\begin{itemize}
  \item \textbf{Stelle} --- kleinste organisatorische Einheit; Aufgabenbündel für eine
    Person.
  \item \textbf{Instanz} --- Stelle mit \emph{Leitungs- und Weisungsbefugnis} gegenüber
    untergeordneten Mitarbeitern (z.\,B.\ Bereichsleitung).
  \item \textbf{Stabstelle} --- unterstützt und berät eine Instanz, hat aber \emph{keine}
    Weisungsbefugnis (z.\,B.\ Öffentlichkeitsarbeit, Qualitätssicherung).
  \item \textbf{Organigramm} --- grafische Darstellung des Stellengefüges.
  \item \textbf{Leitungsspanne} --- Zahl der einer Instanz direkt unterstellten Stellen;
    \textbf{Leitungstiefe} --- Zahl der Hierarchieebenen.
\end{itemize}

\subsection{Leitungssysteme}

\begin{center}
\begin{tabular}{p{2.8cm} p{5.1cm} p{4.6cm}}
\toprule
\textbf{System} & \textbf{Merkmal} & \textbf{Bewertung}\\
\midrule
Einliniensystem & jede Stelle hat nur einen Vorgesetzten & klar, aber lange Dienstwege\\
Mehrliniensystem & eine Stelle hat mehrere Fachvorgesetzte & kurze Wege, aber Kompetenzkonflikte\\
Stabliniensystem & Einliniensystem plus beratende Stabstellen & entlastet die Leitung, aber Stab-Linie-Konflikte\\
Matrixsystem & doppelte Unterstellung nach Funktion und Objekt/Projekt & flexibel, aber Doppelunterstellung\\
\bottomrule
\end{tabular}
\end{center}

\subsection{Ablauforganisation und Geschäftsprozesse}

Während die Aufbauorganisation das \emph{Gerüst} bildet, regelt die
\textbf{Ablauforganisation} die zeitlich-räumlichen \emph{Arbeitsabläufe} --- also, wer
was wann und wo erledigt. Betriebliche Abläufe werden als \emph{Geschäftsprozesse}
beschrieben:

\begin{itemize}
  \item \textbf{Kernprozess} --- wertschöpfender Hauptprozess (z.\,B.\ die Produktion der
    Büromöbel).
  \item \textbf{Unterstützungsprozess} --- flankiert den Kernprozess (z.\,B.\ Buchhaltung,
    Personalwesen).
  \item \textbf{Managementprozess} --- steuert das Unternehmen (Planung, Organisation,
    Kontrolle).
\end{itemize}

\section{Leitbild, Unternehmensziele und Nachhaltigkeit}

Das \textbf{Unternehmensleitbild} beschreibt das Selbstverständnis eines Unternehmens:
seine Vision, seine Werte und Grundsätze im Umgang mit Kunden, Mitarbeitenden und
Gesellschaft. Es gibt Orientierung nach innen und außen.

\textbf{Unternehmensziele} lassen sich unterscheiden in \emph{Sachziele} (was
hergestellt/angeboten wird) und \emph{Formalziele} (übergeordnete Erfolgsziele wie
Gewinn, Wirtschaftlichkeit, Rentabilität). Inhaltlich gliedert man sie in:

\begin{itemize}
  \item \textbf{ökonomische Ziele} --- z.\,B.\ Gewinn, Umsatz, Marktanteil, Liquidität;
  \item \textbf{ökologische Ziele} --- z.\,B.\ Ressourcenschonung, Emissions\-reduktion;
  \item \textbf{soziale Ziele} --- z.\,B.\ faire Arbeitsbedingungen, Ausbildung,
    Arbeits\-sicherheit.
\end{itemize}

\subsection{Die drei Säulen der Nachhaltigkeit}

\begin{merksatz}
\textbf{Nachhaltigkeit} ruht auf drei Säulen: \emph{Ökonomie} (wirtschaftlich tragfähig),
\emph{Ökologie} (umweltverträglich) und \emph{Soziales} (gesellschaftlich verantwortungs\-voll).
Nachhaltiges Handeln bedeutet, alle drei dauerhaft in Einklang zu bringen.
\end{merksatz}

Typische betriebliche Handlungsfelder sind z.\,B.\ Energie- und Materialeffizienz,
Ökostrom, Abfallvermeidung und Recycling (Ökologie), faire Löhne, Gesundheits\-schutz und
Ausbildung (Soziales) sowie langfristige Wirtschaftlichkeit und faire Lieferketten
(Ökonomie).

\section{Zusammenfassung}

\begin{enumerate}
  \item Der \textbf{Kaufmannsbegriff} nach HGB unterscheidet \emph{Istkaufmann} (kraft Gesetzes), \emph{Kannkaufmann} (freiwillig) und \emph{Formkaufmann} (kraft Rechtsform).
  \item Die \textbf{Firma} ist der Handelsname aus Firmenkern und Firmenzusatz; sie muss die fünf Firmengrundsätze (Wahrheit, Klarheit, Ausschließlichkeit, Beständigkeit, Öffentlichkeit) beachten.
  \item Das \textbf{Handelsregister} führt Abteilung A (Einzelkaufleute, Personengesellschaften) und B (Kapitalgesellschaften); Eintragungen wirken \emph{deklaratorisch} oder \emph{konstitutiv}, § 15 HGB regelt die Publizität.
  \item \textbf{Rechtsformen} unterscheiden sich in Haftung und Kapital: Einzelunternehmung/OHG/KG (Personen, Privathaftung) gegenüber GmbH/UG/AG/eG (juristische Personen, beschränkte Haftung; 25.000 €/50.000 € Mindestkapital).
  \item \textbf{Prokura} (umfassend, eingetragen, „ppa.") und \textbf{Handlungsvollmacht} (enger, nicht eingetragen, „i.\,V."/„i.\,A.") sind die kaufmännischen Vollmachten.
  \item \textbf{Aufbauorganisation} (Stelle, Instanz, Stabstelle, Leitungssysteme) und \textbf{Ablauforganisation} (Kern-, Unterstützungs-, Managementprozesse) strukturieren den Betrieb; \textbf{Leitbild, Ziele} und \textbf{Nachhaltigkeit} (drei Säulen) geben die Richtung vor.
\end{enumerate}

\newpage

\section*{Glossar}
\addcontentsline{toc}{section}{Glossar}

\glossareintrag{Istkaufmann}{Wer ein Handelsgewerbe betreibt, das einen kaufmännisch eingerichteten Geschäftsbetrieb erfordert (§ 1 HGB); Eintragung wirkt deklaratorisch.}

\glossareintrag{Kannkaufmann}{Kleingewerbetreibende sowie land-/forstwirtschaftliche Unternehmen, die sich freiwillig ins Handelsregister eintragen lassen können (§§ 2, 3 HGB).}

\glossareintrag{Formkaufmann}{Kapitalgesellschaft (AG, GmbH, KGaA, eG), die kraft ihrer Rechtsform Kaufmann ist; Eintragung wirkt konstitutiv.}

\glossareintrag{Firma}{Der Name, unter dem ein Kaufmann seine Geschäfte betreibt (§ 17 HGB); besteht aus Firmenkern und Firmenzusatz.}

\glossareintrag{Firmengrundsätze}{Firmenwahrheit, Firmenklarheit, Firmenausschließlichkeit, Firmenbeständigkeit und Firmenöffentlichkeit.}

\glossareintrag{Firmenwahrheit}{Grundsatz (§ 18 HGB), dass die Firma nicht über geschäftliche Verhältnisse irreführen darf und Unterscheidungskraft besitzen muss.}

\glossareintrag{Rechtsformzusatz}{Teil des Firmenzusatzes, der die Rechtsform kennzeichnet (z.\,B.\ e.K., OHG, KG, GmbH, AG).}

\glossareintrag{Handelsregister}{Elektronisch geführtes amtliches Verzeichnis beim Amtsgericht; Abteilung A für Personenunternehmen, Abteilung B für Kapitalgesellschaften.}

\glossareintrag{Deklaratorisch}{Rechtsbekundende Eintragung, die eine bereits bestehende Rechtslage nur bestätigt (z.\,B.\ Istkaufmann, Prokura).}

\glossareintrag{Konstitutiv}{Rechtsbegründende Eintragung, mit der die Rechtslage erst entsteht (z.\,B.\ Entstehung einer GmbH oder AG).}

\glossareintrag{Publizität (§ 15 HGB)}{Öffentlicher Glaube des Handelsregisters: Auf eingetragene und bekanntgemachte Tatsachen darf sich jeder Dritte verlassen.}

\glossareintrag{Einzelunternehmung}{Von einer Person geführtes Unternehmen (e.K.); unbeschränkte Haftung mit dem Privatvermögen.}

\glossareintrag{OHG}{Offene Handelsgesellschaft; Gesellschafter haften unbeschränkt, unmittelbar und solidarisch.}

\glossareintrag{Kommanditgesellschaft (KG)}{Personengesellschaft mit unbeschränkt haftendem Komplementär und auf die Einlage beschränkt haftendem Kommanditisten.}

\glossareintrag{GmbH}{Kapitalgesellschaft mit mind.\ 25.000 € Stammkapital; Haftung beschränkt auf das Gesellschaftsvermögen.}

\glossareintrag{UG (haftungsbeschränkt)}{„Mini-GmbH", gründbar ab 1 € Stammkapital, muss Gewinne bis 25.000 € Stammkapital zurücklegen.}

\glossareintrag{Aktiengesellschaft (AG)}{Kapitalgesellschaft mit mind.\ 50.000 € in Aktien zerlegtem Grundkapital; Organe: Hauptversammlung, Vorstand, Aufsichtsrat.}

\glossareintrag{Genossenschaft (eG)}{Der Förderung ihrer Mitglieder verpflichtete juristische Person; Kapital aus Genossenschaftsanteilen.}

\glossareintrag{Prokura}{Umfassende, ins Handelsregister eingetragene Handlungsvollmacht (§§ 48 ff.\ HGB); Unterschrift „ppa.".}

\glossareintrag{Handlungsvollmacht}{Formlose, nicht eingetragene Vollmacht (§ 54 HGB) als General-, Art- oder Einzelvollmacht; Unterschrift „i.\,V." oder „i.\,A.".}

\glossareintrag{Instanz}{Stelle mit Leitungs- und Weisungsbefugnis gegenüber untergeordneten Mitarbeitern.}

\glossareintrag{Stabstelle}{Beratende Stelle ohne Weisungsbefugnis, die eine Instanz unterstützt.}

\glossareintrag{Nachhaltigkeit}{Konzept aus drei Säulen: Ökonomie, Ökologie und Soziales dauerhaft in Einklang bringen.}

\end{document}
